\cotefile{91}
\scanpage{19}
% archive foliation: 18
\unsure{marqués} de $(\unsure{+|4})$, dont les images dans le coin
\unsure{droit du bas} sont les $\unsure{\Longrightarrow|\longrightarrow}$, on cherche un élément
du coin gauche en haut qui donne ces
éléments marqués. On voit qu'il suffit pour
ceci que $P(u_0,S)\ne\varnothing$, car alors le diagramme
$(\unsure{1})$ est isomorphe à $(2)$, or $(\unsure{1|2})$ est
\uline{cartésien}, grâce au fait que $M_{v_0}\longrightarrow M_{u_0}$ est
injectif. Or l'obstruction à trouver un
élément dans $P(u_0,S)$ est \deleted{\unsure{dans}}
\[
\delta(u_0)=\xi\in\operatorname{Ext}^{1}_{\underline{O}_{P_0}}
(P_0;\unsure{I|\mathcal{I}}_{\Gamma_0},\underline{O}_{\Gamma_0}\otimes_{S_0}J)
\]
où $P_0=X_0\times_{S_0}Y_0$, \deleted{\unsure{$\unsure{I|\mathcal{I}}$}} $\Gamma_0$ est le sous-préschéma
de $P_0$ graphe de $u_0$, et où $\unsure{I|\mathcal{I}}_{\Gamma_0}$ est l'Idéal de $P_0$
qui le définit. C'est un module sur $A_0$,
qui s'annule quand on \deleted{\unsure{tensorise par} $\unsure{A_{0\eta}}$}
\inserted{\unsure{localise en} $\unsure{\eta}$}, donc
Si $\pi$ est un paramètre de $A$, il existe un
$n\geq0$ tel que $\pi^{\unsure{n+1|-n}}\xi=0$. Utilisons la suite exacte
\[
(*)\qquad 0\longrightarrow J\xrightarrow{\pi^{\unsure{n+1|-n}}}J\longrightarrow J_n\longrightarrow0
\]
et la suite correspondante
\[
0\longrightarrow\underline{O}_{\Gamma_0}\otimes_{S_0}J
\xrightarrow{\pi^{\unsure{n+1|-n}}}\underline{O}_{\Gamma_0}\otimes_{S_0}J
\longrightarrow\underline{O}_{\Gamma_0}\otimes_{S_0}J_n\longrightarrow0
\]
et la suite exacte des $\operatorname{Ext}^{i}_{\underline{O}_{P_0}}(P_0;\unsure{I|\mathcal{I}}_{\Gamma_0},-)$ : on
trouve que $\xi$ provient d'un élément
\[
\xi'\in\operatorname{Hom}_{\underline{O}_{P_0}}(P_0;\unsure{I|\mathcal{I}}_{\Gamma_0},\underline{O}_{\Gamma_0}\otimes_{S_0}J_n)
\]
\[
\unsure{\bmod}\ \operatorname{Im}\operatorname{Hom}_{\underline{O}_{P_0}}
(P_0;\unsure{I|\mathcal{I}}_{\Gamma_0},\underline{O}_{\Gamma_0}\otimes_{S_0}J).
\]
\margin{\unsure{donc vérifier qu'il ne faut pas remplacer} $P_0$ \unsure{par} $P_1$ \unsure{ci-dessus}.}
